We formalise the LRAM architecture as a composition of typed operators acting on a problem space. Throughout this section, every equation is a standalone display block, suitable for pasting into Microsoft Word Equation Editor.

\subsection{Notation}
Let $\Sigma$ denote the raw syntactic problem space, $\mathcal{A}$ the abstract representation space produced by Prajna, $\mathcal{C} \subseteq \mathcal{A}$ a bounded candidate set, $\mathcal{S} \subseteq \mathcal{C}$ the set of candidates accepted by the oracle, and $\mathcal{V}$ the set of candidates that pass full tier-3 validation.

\subsection{Prajna Abstraction Operator}
The Prajna operator maps a raw problem statement to a compressed abstract representation that preserves the invariants relevant to solution.
\begin{equation}
P : \Sigma \rightarrow \mathcal{A}, \qquad P(\sigma) = \arg\min_{a \in \mathcal{A}} \; \bigl[ \; D(\sigma \,\|\, \mathrm{decode}(a)) + \lambda \, \mathrm{Desc}(a) \; \bigr]
\end{equation}


\subsection{Bounded Candidate Oracle}
The oracle is a predicate that, for each candidate in the bounded set, decides whether the candidate is accepted by a classical efficient check.
\begin{equation}
f : \mathcal{C} \rightarrow \{0, 1\}, \qquad f(c) = 1 \iff c \in \mathcal{S} \subseteq \mathcal{C}
\end{equation}


\subsection{Grover Candidate Amplification}
Given a bounded candidate set of size $N$ and a set of accepted candidates of size $M$, Grover amplitude amplification delivers an accepted candidate after a number of oracle queries proportional to the square root of $N/M$.
\begin{equation}
N_{\mathrm{iter}} \; = \; \left\lfloor \frac{\pi}{4} \sqrt{\frac{N}{M}} \right\rfloor, \qquad N = |\mathcal{C}|, \qquad M = |\mathcal{S}|
\end{equation}


\subsection{Three-Tier Acceptance Probability}
The overall probability that a candidate returned by the Grover stage survives tier-3 validation decomposes multiplicatively into the Prajna faithfulness, the oracle recall, and the validator soundness.
\begin{equation}
\Pr\bigl[c \in \mathcal{V}\bigr] \; = \; \Pr\bigl[c \in \mathcal{A}_{\mathrm{faithful}}\bigr] \cdot \Pr\bigl[c \in \mathcal{S} \,\mid\, c \in \mathcal{A}_{\mathrm{faithful}}\bigr] \cdot \Pr\bigl[c \in \mathcal{V} \,\mid\, c \in \mathcal{S}\bigr]
\end{equation}


\subsection{Do-Attention}
Do-attention replaces conditional attention with interventional attention, computing the key-value score under a structural intervention $\mathrm{do}(x)$ rather than a mere observation $x$.
\begin{equation}
\mathrm{Attn}_{\mathrm{do}}(Q, K, V \,;\, x) \; = \; \mathrm{softmax}\!\left( \frac{Q \, K^{\top}}{\sqrt{d_k}} + \beta \, \Delta(\mathrm{do}(x)) \right) V
\end{equation}


\subsection{World-Model Distillation Loss}
The LRAM world model is distilled from teacher trajectories by a Bellman-consistent objective that combines next-state prediction, reward prediction, and a Kullback-Leibler term aligning the student policy to the teacher ensemble.
\begin{equation}
\mathcal{L}_{\mathrm{WM}}(\theta) \; = \; \mathbb{E}_{(s, a, r, s') \sim \mathcal{D}} \Bigl[ \; \| \hat{s}'_\theta(s, a) - s' \|^2 + \alpha \, (\hat{r}_\theta(s, a) - r)^2 + \gamma \, \mathrm{KL}\!\bigl(\pi_\theta(\cdot \mid s) \,\|\, \bar{\pi}_{\mathrm{teach}}(\cdot \mid s)\bigr) \; \Bigr]
\end{equation}


\subsection{Teacher-Student Distillation Objective}
The student model is trained by minimising a composite loss that combines a Kullback-Leibler term against each teacher in the ensemble, an action-value term against a reference Q-function, and a regularisation term.
\begin{equation}
\mathcal{L}_{\mathrm{Distill}}(\theta) \; = \; \sum_{k=1}^{K} w_k \, \mathrm{KL}\!\bigl(\pi_\theta \,\|\, \pi_{\mathrm{teach},\,k}\bigr) \; + \; \eta \, \mathbb{E}_{(s, a)} \bigl[ (Q_\theta(s, a) - Q_{\mathrm{ref}}(s, a))^2 \bigr] \; + \; \mu \, \Omega(\theta)
\end{equation}


\subsection{Reflexion Update Rule}
The Reflexion outer loop updates an episodic memory by reflecting on failed trajectories and appending textual corrections that modulate future policies.
\begin{equation}
\mathcal{M}_{t+1} \; = \; \mathrm{Decay}(\mathcal{M}_{t}, \rho) \; \cup \; \bigl\{ \; r_{t} = \mathrm{Reflect}(\tau_{t}, y_{t}, v_{t}) \; : \; v_{t} = 0 \; \bigr\}
\end{equation}


\subsection{Convergence Bound for LRAM}
Under the bounded-oracle and faithfulness assumptions, the expected number of oracle queries required to find an accepted and validated candidate is bounded by a closed-form expression.
\begin{equation}
\mathbb{E}\bigl[T_{\mathrm{LRAM}}\bigr] \; \leq \; \frac{1}{\eta_{P} \cdot \eta_{V}} \cdot \frac{\pi}{4} \sqrt{\frac{|\mathcal{C}|}{|\mathcal{S}|}} \; + \; T_{\mathrm{Reflexion}}(\epsilon)
\end{equation}


\subsection{Undecidability Boundary}
We state the formal boundary beyond which no finite proof-search procedure can guarantee termination.
\begin{equation}
\forall \, \mathcal{F} \supseteq \mathrm{PA}, \; \; \exists \, \varphi \in \mathcal{L}(\mathcal{F}) \; : \; \mathcal{F} \nvdash \varphi \; \wedge \; \mathcal{F} \nvdash \neg \varphi \; \wedge \; \mathbb{N} \models \varphi
\end{equation}


\subsection{Summary of the Architecture}
The composition of the operators is a pipeline: a raw problem flows through Prajna into a compressed abstract representation; the compressed representation indexes a bounded candidate set; Grover amplifies the subset accepted by the oracle; tier-3 validates by formal, causal, and circuit checks; Reflexion writes any failure into episodic memory. The inner loop is quantum; the outer loop is classical and reflexive.

\subsection{Full Pipeline Summary}
\begin{equation}
\mathcal{P}
\xrightarrow{8Q}
\text{questioned structure}
\xrightarrow{\text{Prajna}}
\text{compressed forms}
\xrightarrow{\text{TSE}}
\{\mathcal{P}'_1,\mathcal{P}'_2,\ldots\}
\xrightarrow{\text{SSR}+\text{OBM}}
(\mathcal{X},V)
\xrightarrow{\text{Grover/Classical}}
x^{*}
\xrightarrow{\text{Tier-3}}
\text{validated result}
\end{equation}


The complete LRAM pipeline integrates six stages. The 8Q intelligence layer receives the raw problem $\mathcal{P}$ and applies eight dimensions of structured questioning---optimisation, abstraction, constraint encoding, space structuring, execution, measurement, learning, and meta-reflection---to produce a questioned structure: a typed symbolic representation with explicit objects, operations, constraints, and target property. Prajna abstraction then compresses this structure into a family of compact abstract representations $\{\mathcal{P}'_1, \mathcal{P}'_2, \ldots\}$, each preserving the invariants necessary for solution while discarding irrelevant syntactic complexity.

The Transformation Space Engine (TSE) searches the space of problem transformations, evaluating each compressed form against the scoring objective and maintaining a portfolio of candidate worlds. When a world exceeds the handoff threshold, the Search Space Reducer (SSR) constructs the bounded candidate set $\mathcal{X}$ and the Oracle Builder Module (OBM) constructs the verification predicate $V$. This pair $(\mathcal{X}, V)$ is the structured input that enables the next stage.

Grover amplitude amplification (or classical search when quantum hardware is unavailable or when $|\mathcal{X}|$ is small) searches $\mathcal{X}$ using $V$ as the oracle, delivering an accepted candidate $x^*$ in $O(\sqrt{|\mathcal{X}|/M})$ oracle queries where $M$ is the number of solutions. Finally, Tier-3 validation subjects $x^*$ to the full three-layer check---formal proof verification, causal consistency, and circuit-level soundness---before the result is accepted.

The pipeline is self-refining: failures at any stage propagate back through the Reflexion loop, updating episodic memory and biasing future transformation selection and oracle construction. The key architectural insight is that quantum speedup applies only after the problem has been structurally compressed by 8Q, Prajna, and TSE. Attempting Grover search on the raw problem $\mathcal{P}$ without this compression pipeline provides no practical advantage and may provide none at all.
